\documentclass[aspectratio=169, 12pt]{beamer}

% --- Language and Fonts ---
\usepackage{fontspec}
\usepackage[english]{babel}
\usepackage{amsmath, amssymb, bm}
% tikz
\usepackage{tikz}
\usetikzlibrary{shapes, arrows.meta, positioning, shadows, calc}

% --- Bibliography (BibTeX + natbib) ---
\usepackage{natbib}
\bibliographystyle{chicago} 

% =========================================
% ===== keybox ====================
% =========================================
\usepackage[most]{tcolorbox}
\newtcolorbox{keybox}[1]{
  enhanced,
  colback=headbg!4,
  colframe=headbg,
  coltitle=headfg,
  colbacktitle=headbg,
  title=#1,
  fonttitle=\bfseries,
  boxrule=0.8pt,
  arc=1.5mm,
  left=2mm,
  right=2mm,
  top=1mm,
  bottom=1mm,
  before skip=0.2cm,
  after skip=0.2cm
}

\usepackage{booktabs}

% =========================================
% ===== Custom Footline (Page Number) =====
% =========================================


\setbeamercolor{footline}{fg=headbg}
\setbeamerfont{footline}{size=\scriptsize}

\setbeamertemplate{footline}{%
  \hfill%
  \usebeamercolor[fg]{footline}%
  \usebeamerfont{footline}%
  \insertframenumber\,/\,\inserttotalframenumber\kern1em\vskip3pt%
}

\setbeamertemplate{navigation symbols}{}

\usetheme{default}
\usepackage{xcolor}

\definecolor{headbg}{RGB}{30, 60, 120}
\definecolor{headfg}{RGB}{255, 255, 255}
\definecolor{accent}{RGB}{200, 50, 50}

\setbeamercolor{headline}{bg=headbg, fg=headfg}
\setbeamercolor{section in head/foot}{bg=headbg, fg=headfg}
\setbeamercolor{subsection in head/foot}{bg=headbg!80, fg=headfg}
\setbeamercolor{frametitle}{bg=headbg, fg=headfg}

% Step indicator command
\newcommand{\stepindicator}[1]{%
  \hfill {\scriptsize \textit{\textcolor{headbg}{Step #1 / 6}}}%
}
\DeclareMathOperator*{\argmin}{arg\,min}
\newcommand{\indep}{\perp\!\!\!\perp}
\newcommand{\E}{\mathbb{E}}
\newcommand{\Var}{\operatorname{Var}}

\usepackage{hyperref}


% --- Title ---
\title[]{A First Course in Causal Inference Ch.~11}
\author[Kodai Minesaki]{Kodai Minesaki}
\institute[Faculty of Engineering, The University of Tokyo]{Faculty of Engineering, The University of Tokyo}
\date{\today}

\begin{document}

% --- Title Slide ---
\begin{frame}
\titlepage
\end{frame}


\begin{frame}

  {
    \large
    Ch.~12 The Doubly Robust or the Augmented Inverse Propensity Score Weighting Estimator for the Average Causal Effect
  }
  
\end{frame}

\begin{frame}{Ch.~12 The Doubly Robust or the Augmented Inverse Propensity Score Weighting Estimator for the Average Causal Effect}
  \begin{keybox}{Assumption~12.0}
    Ignorability $Z\indep\{Y(1),Y(0)\}\mid X$ and overlap $0<e(X)<1$.
  \end{keybox}
  Under Assum.~12.0, we have two identification formula of the average causal effect $\tau=\E[Y(1)-Y(0)].$ First, the outcome regression formula is
  \begin{equation}
    \tau =\E[\mu_1(X)]-\E[\mu_0(X)],
    \label{eq:outcome-reg}
  \end{equation}
  where,
  \begin{gather*}
    \mu_1(X) = \E[Y(1)\mid X] = \E[Y\mid Z=1,X]\\
    \mu_0(X) = \E[Y(0)\mid X] = \E[Y\mid Z=0,X].
  \end{gather*}
\end{frame}

\begin{frame}
  Second, the IPW formula is
  \begin{equation}
    \tau = \E\left[\frac{ZY}{e(X)}\right] - \E\left[\frac{(1-Z)Y}{1-e(X)}\right],
    \label{eq:ipw}
  \end{equation}
  where $e(X)=\E[Z=1\mid X]$.
  Eq.~\eqref{eq:outcome-reg} is consistent if the outcome model is correctly specified, and Eq.~\eqref{eq:ipw} is consistent if the propensity score model is correctly specified.
  \begin{itemize}
    \item What if one of them is not correctly specified?
    \item We are then motivated to propose an estimator that is consistent if either of them is correctly specified.
    \item It is called the \textit{doubly robust} estimator \citep{scharfstein1999adjusting,bang2005doubly}.
  \end{itemize}
\end{frame}

\begin{frame}{12.1 The doubly robust estimator}
  \begin{itemize}
    \item We posit a model for the conditional means of the outcome $\mu_1(X,\beta_1),\mu_0(X,\beta_0)$, with parameters $\beta_1,\beta_0$.
    \item If the outcome model is corectly specified, then $\mu_1(X)=\mu_1(X,\beta_1)$ and $\mu_0(X)=\mu_0(X,\beta_0)$.
    \item We posit a model for the propensity score $e(X,\alpha)$, with a parameter $\alpha$.
    \item If the propensity score model is correctly specified, then $e(X)=e(X,\alpha)$.
    \item In practice, both models may be misspecified.
    \item Sometimes, we call them \textit{working models} due to the possibility of misspecification.
  \end{itemize}
\end{frame}

\begin{frame}{12.1 The doubly robust estimator}
  Define
  \begin{gather}
    \tilde{\mu}_1^{\text{dr}} = \E\left[\frac{Z\{Y-\mu_1(X,\beta_1)\}}{e(X,\alpha)}+\mu_1(X,\beta_1)\right],\label{eq:dr-outcome-1}\\
    \tilde{\mu}_0^{\text{dr}} = \E\left[\frac{(1-Z)\{Y-\mu_0(X,\beta_0)\}}{1-e(X,\alpha)}+\mu_0(X,\beta_0)\right],\label{eq:dr-outcome-0}
  \end{gather}
  which can be written as
  \begin{gather}
    \tilde{\mu}_1^{\text{dr}} = \E\left[\frac{ZY}{e(X,\alpha)}-\frac{Z-e(X,\alpha)}{e(X,\alpha)}\mu_1(X,\beta_1)\right],\label{eq:dr-ipw-1}\\
    \tilde{\mu}_0^{\text{dr}} = \E\left[\frac{(1-Z)Y}{1-e(X,\alpha)}-\frac{e(X,\alpha)-Z}{1-e(X,\alpha)}\mu(X,\beta_0)\right].\label{eq:dr-ipw-0}
  \end{gather}
\end{frame}

\begin{frame}{12.1 The doubly robust estimator}
  \begin{itemize}
    \item Eq.~\eqref{eq:dr-outcome-1},~\eqref{eq:dr-outcome-0} augment the outcome regression estimator by inverse propensity score weighting terms of the residuals ($Y-\mu_Z(X,\beta_Z)$).
    \item Eq.~\eqref{eq:dr-ipw-1},~\eqref{eq:dr-ipw-0} augment the IPW estimator by the imputed outcomes.
    \item For this reason, the doubly robust estimator is also called the \textit{augmented inverse propensity score weighting} (AIPW) estimator.
    \item The augmentation strengthens the theoretical properties in the following sense.
  \end{itemize}
  \begin{keybox}{Theorem~12.1}
    Assume ignorability $Z\indep \{Y(1),Y(0)\}\mid X$ and overlap $0<e(X)<1$. 
    \begin{enumerate}
      \item If either $e(X,\alpha)=e(X)$ or $\mu_1(X,\beta_1)=\mu_1(X)$, then $\tilde{\mu}_1^{\text{dr}}=\E[Y(1)]$.
      \item If either $e(X,\alpha)=e(X)$ or $\mu_0(X,\beta_1)=\mu_0(X)$, then $\tilde{\mu}_0^{\text{dr}}=\E[Y(0)]$.
      \item If either $e(X,\alpha)=e(X)$ or $\{\mu_1(X,\beta_1)=\mu_1(X),\mu_0(X,\beta_1)=\mu_0(X)\}$, then $\tilde{\mu}_1^{\text{dr}} - \tilde{\mu}_0^{\text{dr}} = \tau.$
    \end{enumerate}
  \end{keybox}
\end{frame}

\begin{frame}{12.1 The doubly robust estimator}
  \begin{itemize}
    \item By Thm.~12.1, \textbf{if either the propensity score model or the outcome model is correctly specified,} $\tilde{\mu}_1^{\text{dr}} - \tilde{\mu}_0^{\text{dr}} = \tau.$
  \end{itemize}
  \small
  \begin{proof}
  \vspace{-3mm}
    \begin{align*}
      \tilde{\mu}_1^{\text{dr}} - \E[Y(1)]
      &= \E\left[\frac{Z\{Y(1)-\mu_1(X,\beta_1)\}}{e(X,\alpha)}+\mu_1(X,\beta_1) - Y(1)\right]\\
      &= \E\left[\frac{Z-e(X,\alpha)}{e(X,\alpha)}(Y(1)-\mu_1(X,\beta_1))\right]\\
      &= \E\left[\E\left[\frac{Z-e(X,\alpha)}{e(X,\alpha)}(Y(1)-\mu_1(X,\beta_1))\mid X\right]\right]\quad (\because \text{LIE})\\
      &= \E\left[\E\left[\frac{Z-e(X,\alpha)}{e(X,\alpha)}\mid X\right]\times \E\left[Y(1)-\mu_1(X,\beta_1)\mid X\right]\right]\quad (\because\text{strong ignorability})\\
      &= \E\left[\frac{e(X)-e(X,\alpha)}{e(X,\alpha)}\times \{\mu_1(X)-\mu_1(X,\alpha)\}\right].
    \end{align*}
  \end{proof}
\end{frame}

\begin{frame}{12.1 The doubly robust estimator}
  We can obtain their sample analogs to construct a doubly robust estimator for $\tau$.
  
  \begin{keybox}{12.1~Definition (doubly robust estimator for the average causal effect)}
    Based on the data $(X_1,Z_i,Y_i)_{i=1}^n$, we can obtain a doubly robust estimator for $\tau$ by the following steps:
    \begin{enumerate}
      \item obtain the fitted values of the propensity scores: $e(X_i,\hat{\alpha})$;
      \item obtain the fitted values of the outcome means: $\mu_1(X_i,\hat{\beta}_1),\mu_0(X_i,\hat{\beta}_0)$;
      \item construct the doubly robust estimator: $\hat{\tau}^{dr}=\hat{\mu}_1^{\text{dr}}-\hat{\mu}_0^{\text{dr}}$, where
      \begin{gather*}
        \hat{\mu}_1^{\text{dr}} = \frac{1}{n}\sum_{i=1}^{n}\left[\frac{Z_i\{Y_i-\mu_1(X_i,\hat{\beta}_1)\}}{e(X_i,\hat{\alpha})} + \mu_1(X_i,\hat{\beta}_1)\right],\\
        \hat{\mu}_0^{\text{dr}} = \frac{1}{n}\sum_{i=1}^{n}\left[\frac{(1-Z_i)\{Y_i-\mu_0(X_i,\hat{\beta}_0)\}}{1-e(X_i,\hat{\alpha})} + \mu_0(X_i,\hat{\beta}_0)\right]
      \end{gather*}
    \end{enumerate}
  \end{keybox}
\end{frame}

\begin{frame}{12.1 The doubly robust estimator}
  By definition, we can also write the doubly robust estimator as:
  \[
  \hat{\tau}^{\text{dr}} = \hat{\tau}^{\text{reg}} + \frac{1}{n}\sum_{i=1}^{n}\frac{Z_i\{Y_i-\mu_1(X_i,\hat{\beta}_1)\}}{e(X,\hat{\alpha})} - \frac{1}{n}\sum_{i=1}^{n}\frac{(1-Z_i)\{Y_i-\mu_0(X_i,\hat{\beta}_0)\}}{1-e(X_i,\hat{\alpha})}.
  \]
  We can also rewrite it as
  \[
  \hat{\tau}^{\text{dr}} = \hat{\tau}^{\text{ipw}} - \frac{1}{n}\sum_{i=1}^{n}\frac{Z_i-e(X_i,\hat{\alpha})}{e(X_i,\hat{\alpha})}\mu_1(X_i,\hat{\beta}_1) + \frac{1}{n}\sum_{i=1}^{n}\frac{e(X_i,\hat{\alpha})-Z_i}{1-e(X_i,\hat{\alpha})}\mu_0(X_i,\hat{\beta}_0).
  \]
\end{frame}

\begin{frame}{12.2 More intuition and theory for the doubly robust estimator}
  The IPW estimator for $\mu_1$ based on $\mu_1 = \E\left[ZY/e(X)\right]$ completely ignores the outcome model of $Y$. With possibly misspecified outcome model $\mu_1(X,\beta_1)$,
  \[
  \mu_1 =\E[Y(1)] = \E[Y(1)-\mu_1(X,\beta_1)] + \E[\mu_1(X,\beta_1)].
  \]
  If we apply the IPW formula to the first term in the above formula:
  \[
  \mu_1= \E\left[\frac{Z(Y(1)-\mu_1(X,\beta_1))}{e(X)}\right] + \E[\mu_1(X,\beta_1)] = \E\left[\frac{Z(Y(1)-\mu_1(X,\beta_1))}{e(X)}+\mu_1(X,\beta_1)\right],
  \]
  which holds if the propensity core model correct without assuming that the outcome model is correct. 
\end{frame}

\begin{frame}{Future Reading}
  \begin{itemize}
    \item Chernozhukov, V., Chetverikov, D., Demirer, M., Duflo, E., Hansen, C., Newey, W., \& Robins, J. (2018). Double/debiased machine learning for treatment and structural parameters. \url{https://academic.oup.com/ectj/article/21/1/C1/5056401}
  \end{itemize}
\end{frame}

% =========================

% =========================
\section{}
\begin{frame}[allowframebreaks, noframenumbering]{References}
  \beamerdefaultoverlayspecification{}
  \small
  \bibliography{../references}
\end{frame}

\end{document}